Студент выбирает \textbf{один вариант} и сдаёт только его. Вариант 1 ---
базовая программа факториала на C; вариант 2 --- описанная ниже
усложнённая работа с BIOS и QEMU. Порядок показа,
самостоятельной работы и защиты одинаков; не требуйте заданий или
отчёта из другого варианта.

Для варианта 1 проверьте собственный \texttt{factorial.c}, сборку с
\texttt{-g\ -O0}, запуск на 0, 5, 20 и ошибочном вводе, ненулевой код
выхода при ошибке, раздельные \texttt{stdout}/\texttt{stderr} и
объяснение работы цикла в GDB. Попросите найти ошибку
\texttt{i\ \textless{}\ n} (не умножается на \texttt{n}) и исправить её.
На проверку студент приносит свой исходник и отчёт, а не только
исполняемый файл. Эталон для локальных проверок
\texttt{factorial/reference/factorial.c} не включается в студенческий
архив.

Ниже --- подсказки и ответы только для варианта 2.

Цель: студенты на Debian вручную устанавливают инструменты, собирают
свой BIOS-сектор и полный журнал, запускают QEMU, проверяют результат по
экрану, COM1 и байтам образа и записывают наблюдения в отчёте.
Преподаватель не предлагает студенту \texttt{make} или автоматическую
проверку. Рекомендуемая среда --- Debian Linux актуальной версии на
компьютере или в виртуальной машине. Можно работать и в другом
дистрибутиве Linux с теми же инструментами; инструкции по установке
пакетов написаны для Debian. При выборе другой системы студент
самостоятельно решает вопросы с её пакетами и настройкой: преподаватели
не могут изучить каждый дистрибутив отдельно. Для домашнего опыта с
C/GDB нужна среда Linux x86-64. Если среда ещё не готова, студент может
потратить семинар на установку системы или настройку виртуальной машины,
показать достигнутый результат в конце занятия и завершить оставшиеся
задания дома. Пакеты студент устанавливает сам. Следующий этап
начинается после проверки предыдущего.

\begin{longtable}[]{@{}p{.19\linewidth}p{.27\linewidth}p{.43\linewidth}@{}}
\toprule
Этап & Наблюдение & При затруднении покажите место, не число \\
\midrule
\endhead
Загрузочный сектор & \texttt{HELLO\ BIOS} после сборки 512-байтного
сектора & \texttt{message}, \texttt{bits\ 16}, \texttt{times},
\texttt{dw} \\
Готовый журнал & журнал без изменений: \texttt{SAVED} → \texttt{LOADED}
& \texttt{boot16.asm}, \texttt{journal.asm}, два разных образа \\
Вывод & новый цвет и позиция заголовка & \texttt{vga\_text} и
\texttt{redraw} перед \texttt{title\_text} \\
Ввод & Q/E вместо S/L & \texttt{handle\_scan\ .normal},
\texttt{help\_text} \\
Вычисления & A=12, оба средних, загрузка 12 &
\texttt{handle\_scan\ .ten}, \texttt{load\_grades\ .check} \\
Диск & \texttt{GR16} в новом секторе копии HDD & \texttt{dap\ dq},
размер сектора 512 \\
Итог & путь клавиша → RAM → VGA → HDD & студенческий отчёт и собственные
наблюдения \\
\bottomrule
\end{longtable}

\hypertarget{ux43fux440ux438ux451ux43c-ux440ux430ux431ux43eux442ux44b}{%
\subsubsection{Приём
работы}\label{ux43fux440ux438ux451ux43c-ux440ux430ux431ux43eux442ux44b}}

В конце семинара №4 попросите студента показать достигнутый результат:
установку Debian, запущенную программу или другую выполненную часть
работы, а также зафиксированные шаги и наблюдения. Предложите задать
вопросы, возникшие при выполнении работы. Одинаковой обязательной стадии
для показа нет. К семинару №5 студент завершает задания, разбирает
использованные команды, понятия и работу программы, готовит отчёт и
сохраняет исходники и образы для повторного запуска. Пока студенты
выполняют семинар №5, принимайте
предыдущую работу индивидуально: отчёт, работающий результат, объяснение
расчётов и ответы на вопросы. Сообщите, если нужно что-то исправить.
На последнем занятии примите защиту предыдущей лабораторной работы и
оставшиеся задолженности, объявите текущую накопленную оценку за
лабораторные работы.

\hypertarget{ux43eux442ux432ux435ux442ux44b-ux43dux430-ux447ux438ux441ux43bux43eux432ux44bux435-ux437ux430ux434ux430ux43dux438ux44f-ux43dux435-ux434ux438ux43aux442ux443ux439ux442ux435-ux434ux43e-ux440ux430ux441ux447ux451ux442ux430}{%
\subsubsection{Ответы на числовые задания (не диктуйте до
расчёта)}\label{ux43eux442ux432ux435ux442ux44b-ux43dux430-ux447ux438ux441ux43bux43eux432ux44bux435-ux437ux430ux434ux430ux43dux438ux44f-ux43dux435-ux434ux438ux43aux442ux443ux439ux442ux435-ux434ux43e-ux440ux430ux441ux447ux451ux442ux430}}

\begin{itemize}
\tightlist
\item
  Цвет текста жёлтый \texttt{Eh}, фон синий \texttt{1h}, бит мигания 0:
  \texttt{(1\ \textless{}\textless{}\ 4)\ \textbar{}\ E\ =\ 1Eh} ---
  новый операнд \texttt{mov\ ah} в \texttt{vga\_text}.
  \texttt{mov\ ax,0x0720} в \texttt{redraw} только очищает экран
  пробелами с исходным атрибутом и не задаёт цвет вновь выводимых
  символов.
\item
  Строка 2, столбец 5: \texttt{BH=02h}, \texttt{BL=05h},
  \texttt{BX=0205h} перед \texttt{title\_text}. Смещение VGA
  \texttt{2×(80×2+5)=330=0x14A} байтов, физический адрес
  \texttt{0xB814A}. Надпись умещается до строки таблицы.
\item
  Старые скан-коды S=31₁₀=\texttt{1Fh}, L=38₁₀=\texttt{26h}; новые
  Q=16₁₀=\texttt{10h}, E=18₁₀=\texttt{12h}. A=30₁₀=\texttt{1Eh} остаётся
  вводом максимальной оценки. Студенту дана таблица десятичных значений,
  а не готовые операнды для NASM.
\item
  В \texttt{.ten} новое значение \texttt{AL=12}; в
  \texttt{load\_grades\ .check} разрешить число 12, иначе сохранение
  получится, но загрузка отвергнет сектор. \texttt{{[}12,0{]}} даёт
  среднее 6.00; \texttt{{[}12,1,1{]}} даёт сумму 14, количество 3 и
  \texttt{(1400+1)//3=467} сотых → 4.67. Делитель 10 в \texttt{vga\_two}
  относится к десятичной записи и не меняется.
\item
  Смещение файла 4096 / размер сектора 512 = LBA 8. В DAP нужно
  \texttt{dq\ 8}; байты сигнатуры \texttt{GR16} начинаются с 4096,
  оценки --- с 4104. Первый запуск со старой копией диска и новым LBA
  обычно выдаёт \texttt{EMPTY}.
\item
  Домашние примеры: ярко-зелёный на чёрном --- \texttt{0Ah}; белый на
  красном --- \texttt{4Fh}. Для домашнего опыта с повреждением первой
  оценки 12 запись 1 по адресу \texttt{OFFSET+8} без коррекции checksum
  должна дать \texttt{BAD\_CHECKSUM}.
\end{itemize}

Полный пример констант: \texttt{solutions/lab/journal.asm}. Он удобен
для подготовки демонстрации, но выдавать вычисленные операнды до попытки
студента не нужно. Исторические заготовки по этапам остаются рядом как
материалы лекции; студент выполняет текущую лабораторную на
\textbf{полном} файле этапа 6, скопированном в рабочий каталог.

\hypertarget{ux447ux442ux43e-ux43fux440ux43eux432ux435ux440ux44fux442ux44c-ux432-ux43eux442ux447ux451ux442ux435}{%
\subsubsection{Что проверять в
отчёте}\label{ux447ux442ux43e-ux43fux440ux43eux432ux435ux440ux44fux442ux44c-ux432-ux43eux442ux447ux451ux442ux435}}

При пересборке студент повторяет NASM для журнала и все три команды
\texttt{dd}: первая создаёт чистую дискету, две другие кладут загрузчик
и код. Диск оценок остаётся прежним, поэтому сохранение можно проверить
после нового запуска.

Для каждого изменения нужны: место в коде, исходное число, формула,
новое число с основанием системы счисления, команды сборки и запуска,
ожидаемый и наблюдаемый результат, объяснение несовпадения при ошибке.
Отдельно спросите, почему загрузчик читает не весь HDD, почему два
\texttt{-drive} не являются двумя именами файла в гостевой ОС и почему
пересборка дискеты не должна затронуть оценки на другом образе.

Типичные ошибки: считать ASCII вместо скан-кода; путать \texttt{BH} и
\texttt{BL}; использовать код жёлтого цвета как фон 0--7; поменять
максимум 12, но не изменить проверку загружаемых байтов; заново создать
\texttt{grades.img} перед проверкой перезапуска; записать бинарный
журнал в первый сектор; испортить основной файл вместо третьей копии.
Первым исправляйте модель адресации, затем числовой операнд.

Программа \texttt{memory.c} дома собирается GCC для Linux x86-64 LP64.
Переход к теме №4: \texttt{file\ memory} узнаёт ELF-файл (точная строка
зависит от версии GCC), \texttt{./memory} запускает процесс, а
\texttt{echo\ \$?} сразу после его завершения показывает \texttt{0},
если обе проверки в C прошли. Не называйте файл на диске процессом:
процесс существует во время выполнения программы. Адрес \texttt{\&x} в
этом пользовательском процессе виртуальный; его нельзя отождествлять с
физическим адресом BIOS-журнала в QEMU. В GDB студент сам вводит
\texttt{break}, \texttt{run}, \texttt{print}, \texttt{x/4xb},
\texttt{x/3dw} и объясняет вывод. Спросите, зачем для \texttt{memory}
нужны услуги Linux и почему \texttt{hello.bin} загружает BIOS, а не
Linux.
